\scanpage{60}
Cela équivaut, en termes de $(\mathcal M,\varphi)$, à la condition que l'on ait

$4^\circ$ a) Dans $\mathcal M$ les sommes disjointes finies existent, et $\varphi$ y commute\deleted{\unsure{nt}}.

b) $\exists$ monomorphisme $e_{\mathcal M}\amalg e_{\mathcal M}\longrightarrow I_0$.

\textcircled{D} Supposons que dans $\mathcal K$ les sommes amalgamées $X\amalg_Y Z$ existent \uline{si} $Y\to X$ et $Y-Z$ mono.
(Alors pour un morph. quelconque $Y\mathrel{\overset{f}{-}}X$, $X\amalg_Y X$ existe, car on peut remplacer $f$ par l'inclusion de son image).
\inserted{$Y\xrightarrow{f}X$}
un épi ssi $X\amalg_Y X\xrightarrow{\nabla_{X/Y}}X$ est iso. Donc si dans $\mathcal M$ les $=$ types de sommes amalgamées existent, et $\varphi$ y commute, alors on a $\varphi(f)$ est un épi ssi $f$ est épi; si dans $\mathcal K$ épi $\Rightarrow$ épi eff., kif kif dans $\mathcal M$; Si dans $\mathcal K$ tout mono est effectif, de $=$ dans $\mathcal M$.

On considère un diagramme dans $\mathcal M$, i.e. d'objets modérés de $\mathcal K$, \inserted{et morphismes modérés}
\diagram{\begin{tikzcd}[ampersand replacement=\&]
\&Y\arrow[dl,"{i}"']\arrow[dr,"{j}"]\&\\
X\&\&Z
\end{tikzcd}}
avec $i,j$ mono, on cherche $X\amalg_Y Z$ \uline{dans} $\mathcal M$ avec $\varphi$ y commutant.
Donc on cherche une structure modérée sur $X\amalg_Y Z$ dans $\mathcal K$, \deleted{et un diag} tel que
$\begin{array}{c}X\longrightarrow X\amalg_Y Z\\ Z\longrightarrow\end{array}$
soient modérés, et que ça satisfasse le Pb universel.

\diagram{\begin{tikzcd}[ampersand replacement=\&]
\&Y\arrow[dl]\arrow[dr]\&\&\\
X\arrow[dr]\&\&Z\arrow[dl]\&\\
\&X\amalg_Y Z\&X\amalg Z\arrow[l]\&Y\arrow[l,shift left]\arrow[l,shift right]
\end{tikzcd}}

\margin{61}
